\subsection{同一个系数为什么既是力除以长度，又是功除以面积}
设薄液膜有一条可滑动的直边，长度为 $\ell$，准静态向外移动 $\dd x$。液膜有前后两个界面，面积总增量为 $\dd A=2\ell\dd x$，外力做功为 $F\dd x$。在温度、组成不变，且过程可逆的条件下，$F\dd x=\gamma\dd A$，所以 $F=2\gamma\ell$。若只有一个界面，相应的因子2消失。
因此 $\mathrm{N/m}=\mathrm{J/m^2}$ 并非偶然的单位巧合，而是力做功与面积变化之间的关系。液膜沿界面拉紧边界，不是处处沿固定空间方向施力。

\subsection{球面压差与毛细升降的推导}
先看单界面的球形液滴。令半径准静态增加 $\dd R$，体积和面积的变化为
\[
 \dd V=4\pi R^2\dd R,\qquad \dd A=8\pi R\dd R.
\]
内外压强差做功为 $(p_{\rm in}-p_{\rm out})\dd V$，用于增加界面自由能 $\gamma\dd A$。约去公共因子后，$p_{\rm in}-p_{\rm out}=2\gamma/R$。薄壁肥皂泡有两个近似同半径的界面，面积变化加倍，压差才成为 $4\gamma/R$；“两个界面”不是把半径也乘2。

再看半径为 $r$ 的竖直毛细管，管外液面宽广。接触角 $\theta$ 从液体内部量取，圆形接触线长 $2\pi r$，界面张力的竖直合力为 $2\pi r\gamma\cos\theta$。在弯月面体积相对液柱体积可忽略、气体密度可忽略的近似下，它平衡高出外液面的液柱重量：
\[
 2\pi r\gamma\cos\theta=\rho g\pi r^2h.
\]
于是得到前式。$h$ 是相对管外液面的有符号高度：润湿时 $\theta<\pi/2$、$h>0$；不润湿时 $\theta>\pi/2$、$h<0$，表示液面下降。忽略弯月面体积通常要求管半径远小于毛细长度 $\sqrt{\gamma/(\rho g)}$。接触角变化、管壁污染或液柱尚在运动时，不能原样套用静态结果。

\begin{origexamplebox}{液滴合并时的界面能变化}\label{ch13:ex:teaching-droplets}
两滴同种不可压缩液体均为半径 $R$ 的球形液滴，合并后形成一个球形液滴。设表面张力为常量 $\gamma$，忽略重力。求合并后的半径、界面能减少量及新液滴的内外压差。界面能的减少是否等于总能量的消失？
\end{origexamplebox}
\begin{xesolution}
先守恒的是体积而不是表面积。设新半径为 $R_f$，则
\[
 \frac43\pi R_f^3=2\frac43\pi R^3,
 \qquad R_f=2^{1/3}R.
\]
初、末总表面积分别为 $8\pi R^2$ 和 $4\pi 2^{2/3}R^2$。因此界面能减少
\[
 \Delta E_s=4\pi\gamma R^2(2-2^{2/3})>0.
\]
这部分能量可先转化为液滴内部流动、形变振动的能量，再由黏性耗散转化为内能；没有凭空消失。若过程中温度明显改变，常量 $\gamma$ 的近似还须重查。
新液滴压差为 $\Delta p_f=2\gamma/(2^{1/3}R)$，比原压差小。半径增大、曲率减小与压差下降一致；$\gamma R^2$ 的量纲是能量，$\gamma/R$ 的量纲是压强。
\end{xesolution}
